\section{Cohomological Classification}

We interpret the transport cocycle as a cohomological object on
$\Gfifteen$. Cohomology of graphs with $\mathbb{Z}_2$ coefficients is
standard (see \cite{hatcher,godsil_royle}).

\subsection{Cochains}

Let $C^1(\Gfifteen,\ZZtwo)$ denote the space of $\ZZtwo$-valued functions
on edges. Elements of $C^1$ assign a value in $\{0,1\}$, equivalently a
sign in $\{\pm 1\}$, to each edge.

Let $C^0(\Gfifteen,\ZZtwo)$ denote the space of $\ZZtwo$-valued functions
on vertices. The coboundary map
\[
\delta : C^0(\Gfifteen,\ZZtwo)\to C^1(\Gfifteen,\ZZtwo)
\]
encodes vertex switching, corresponding to the standard notion of
switching equivalence in signed graphs \cite{zaslavsky}.

\subsection{Classifying cocycle of the transport cover}

The lifted transport $\Gthirty\to\Gfifteen$ defines a function
\[
\varepsilon:E(\Gfifteen)\to\ZZtwo,
\]
where $\varepsilon(e)$ records whether transport across $e$ preserves or
swaps sheet.

\begin{proposition}
The function $\varepsilon$ is the classifying cocycle of the regular
$\ZZtwo$-cover
\[
\Gthirty\to\Gfifteen.
\]
\end{proposition}

\begin{proof}
A regular $\ZZtwo$-cover of a connected graph is determined, up to
equivalence, by a $\ZZtwo$-valued $1$-cocycle recording whether each
base edge lifts sheet-preservingly or sheet-swappingly.

By construction, $\varepsilon(e)$ records exactly this dichotomy for the
cover $\Gthirty\to\Gfifteen$. Hence $\varepsilon$ is the classifying cocycle
of the cover.
\end{proof}

For any cycle $C\subset \Gfifteen$, define the holonomy
\[
\omega(C)=\sum_{e\in C}\varepsilon(e)\pmod 2.
\]

Then $\omega(C)$ computes the sheet parity of the lifted transport around
$C$. In particular, $\varepsilon$ is a $1$-cocycle.

\subsection{Gauge equivalence}

Switching at a vertex adds a coboundary. Hence two edge functions related
by switching represent the same cohomology class. The transport cocycle
therefore determines a well-defined element
\[
[\varepsilon]\in H^1(\Gfifteen,\ZZtwo).
\]

\subsection{Dimension of the cohomology group}

Because $\Gfifteen$ is connected,
\[
\dim H^1(\Gfifteen,\ZZtwo)=|E(\Gfifteen)|-|V(\Gfifteen)|+1.
\]
Since
\[
|V(\Gfifteen)|=15,\qquad |E(\Gfifteen)|=30,
\]
it follows that
\[
\dim H^1(\Gfifteen,\ZZtwo)=30-15+1=16.
\]

Thus the full cohomology group is
\[
H^1(\Gfifteen,\ZZtwo)\cong \ZZtwo^{16}.
\]

The Thalean cocycle is therefore one distinguished nonzero class inside a
$16$-dimensional mod-$2$ cohomology space.

\subsection{Nontriviality}

\begin{theorem}
The transport cocycle $\varepsilon$ represents a nontrivial class in
$H^1(\Gfifteen,\ZZtwo)$.
\end{theorem}

\begin{proof}
A cocycle is trivial if and only if all cycle holonomies vanish.

By the transport construction, there exist cycles in $\Gfifteen$ with
both trivial and nontrivial holonomy. In particular, there exists a cycle
$C$ with $\omega(C)=1$. Hence $\varepsilon$ is not a coboundary and so
defines a nontrivial cohomology class.
\end{proof}

\subsection{Minimal support in the switching class}

Within its switching class, the cocycle $\varepsilon$ admits
representatives with varying support sizes. Define the support
\[
S=\{e\in E(\Gfifteen):\varepsilon(e)=1\}.
\]

\begin{proposition}
The minimal support size of $\varepsilon$ over all switching-equivalent
representatives is $6$.
\end{proposition}

\begin{proof}
An exhaustive search over all vertex switchings, with one vertex fixed to
remove global symmetry, yields representatives with exactly $6$ nonzero
edges. No switching configuration produces fewer.
\end{proof}

Moreover, this minimal representative is not unique: the same exhaustive
search yields $40$ distinct supports of size $6$.

Thus the number $6$ measures the sparsest realization of the Thalean
class, not the dimension of the ambient cohomology group.

\subsection{Interpretation}

The cohomology class $[\varepsilon]$ is the obstruction to trivializing
the $\ZZtwo$-cover $\Gthirty\to\Gfifteen$.

Equivalently, it records the existence of cycles whose lifted transport
returns on the opposite sheet. Its minimal support gives a quantitative
measure of how sparsely this obstruction can be represented after gauge
reduction.
